\documentclass{article}
\usepackage[utf8]{inputenc}
\usepackage{enumitem}

\title{Exemples d'application de la règle des quatres yeux }
\author{Takougang Tamafo Mendel}
\date{22 mars 2026}

\begin{document}

\maketitle

\section*{Définition}
La règle des quatre yeux stipule que toute opération sensible ou décision critique doit être validée par au moins deux personnes indépendantes, afin de réduire les risques d'erreur ou de fraude.

\section*{Exemple 1 : Validation des paiements}
\begin{itemize}[label=--]
    \item \textbf{Situation :} Une entreprise doit effectuer un virement bancaire de 500\,000 Fcfa .
    \item \textbf{Application de la règle :} 
    \begin{enumerate}
        \item Le responsable financier prépare le virement et saisit les informations dans le système.
        \item Un second cadre (contrôleur ou directeur) vérifie les coordonnées bancaires, le montant et la justification avant d’approuver.
    \end{enumerate}
    \item \textbf{Résultat :} Le paiement n’est exécuté que si deux personnes distinctes ont validé l’opération.
\end{itemize}

\section*{Exemple 2 : Contrôle des écritures comptables}
\begin{itemize}[label=--]
    \item \textbf{Situation :} Un auditeur interne enregistre une écriture de correction de 10\,000 Fcfa dans les comptes.
    \item \textbf{Application de la règle :}
    \begin{enumerate}
        \item L’auditeur saisit l’écriture et documente la justification.
        \item Un superviseur ou chef de mission relit et valide l’écriture avant son intégration définitive.
    \end{enumerate}
    \item \textbf{Résultat :} La correction comptable est fiable car elle a été contrôlée par deux personnes indépendantes.
\end{itemize}

\end{document}
